%% Allures des courbes du MOUVEMENT UNIFORMÉMENT ACCÉLÉRÉ :
%% position parabolique, vitesse affine croissante, accélération constante positive.
\documentclass[tikz,border=4pt]{standalone}
\usepackage{liaisons}
\begin{document}
\begin{tikzpicture}[x=1cm,y=1cm, scale=0.92,
  axe/.style={-{Stealth[length=5pt]}, line width=0.6pt},
  courbe/.style={solideA, line width=1.3pt, line cap=round}]

%% \petitrepere{décalage}{grandeur en ordonnée}{tracé}
\newcommand\petitrepere[3]{%
\begin{scope}[shift={(#1,0)}]
  \draw[axe] (-0.18,0) -- (2.15,0) node[anchor=north east, inner sep=3pt, font=\small] {$t$};
  \draw[axe] (0,-0.18) -- (0,1.62);
  \node[anchor=south west, inner sep=1pt, font=\small] at (-0.1,1.62) {#2};
  #3
\end{scope}}

\node[font=\small\bfseries] at (3.6,2.35) {Uniform\'ement acc\'el\'er\'e};

% (1) position : parabole à concavité vers le haut
\petitrepere{0}{$x$ ou $\theta$}{%
  \draw[courbe] plot[domain=0:1.95, samples=40, smooth] (\x,{0.375*\x*\x});}

% (2) vitesse : droite oblique croissante
\petitrepere{2.7}{$v$ ou $\theta'$}{%
  \draw[courbe] (0,0.1) -- (1.95,1.42);}

% (3) accélération : constante positive
\petitrepere{5.4}{$a$ ou $\theta''$}{%
  \draw[courbe] (0,0.85) -- (1.95,0.85);
  \node[font=\small, text=solideA, anchor=south] at (1.0,0.95) {$a > 0$};}
\end{tikzpicture}
\end{document}
